\centerline{\bf \Huge Посчитаем дважды}

\begin{flushright}
        \scriptsize{Не бывает безвьходных ситуаций.\\ Бывают ситуации, выход из которых тебя не устраивает.\\Нара Шикимару}
\end{flushright}

\problem{1}
Можно ли в клетки прямоугольной таблицы записать числа так, что 

а) в каждом столбце сумма положительная, а в каждой строке - отрицательная;

б) в каждом столбце сумма больше 10 , а в каждой строке - меньше 10 ?

\problem{2}
Можно ли записать натуральные числа в клетки таблицы $7 \times 7$ так, чтобы в любом квадрате $2 \times 2$ и любом квадрате $3 \times 3$ сумма чисел была нечётная?

\problem{3}
Можно ли расставить целые числа от 2 до 17 в клетках таблицы $4 \times 4$ так, чтобы суммы во всех строках были равны и не было такой строки, в которой какое-то число делилось бы на другое?

\problem{4}
В клетках таблицы $10 \times 10$ расставлены числа от 1 до 100 . Пусть $S_1, S_2, \ldots, S_{10}$ - суммы чисел, стоящих в столбцах таблицы. Могло ли оказаться так, что в ряду чисел $S_1, S_2, \ldots, S_{10}$ каждые два соседних различаются на 1 ?

\problem{5}
Назовём таблицу интересной, если суммы чисел во всех её строках разные, а суммы чисел во всех её столбцах одинаковые. Существует ли интересная таблица $8 \times 8$ в клетках которой стоят только из нули и единицы?


\problem{6}
Можно ли в таблице $7 \times 9$ закрасить 24 клетки так, чтобы в любом квадрате $3 \times 3$ было ровно 4 закрашенные клетки, а в любом квадрате $5 \times 5$ было не менее 14 незакрашенных клеток?